%%% تنظیمات بستهٔ glossaries برای واژه‌نامه‌ها و فهرست اختصارات

%%% واژه‌نامهٔ فارسی به انگلیسی
\newglossarystyle{myFaToEn}{%
  \renewenvironment{theglossary}{}{}
  \renewcommand*{\glsgroupskip}{\vskip 7mm}
  \renewcommand*{\glsgroupheading}[1]{\subsection*{\glsgetgrouptitle{##1}}}
  \renewcommand*{\glossentry}[2]{%
    \noindent\glsentryname{##1}\dotfill\space\glsentrytext{##1}\par
  }
}

%%% واژه‌نامهٔ انگلیسی به فارسی
\newglossarystyle{myEntoFa}{%
  \renewenvironment{theglossary}{}{}
  \renewcommand*{\glsgroupskip}{\vskip 7mm}
  \renewcommand*{\glsgroupheading}[1]{%
    \begin{LTR}\subsection*{\glsgetgrouptitle{##1}}\end{LTR}}
  \renewcommand*{\glossentry}[2]{%
    \noindent\glsentrytext{##1}\dotfill\space\glsentryname{##1}\par
  }
}

%%% فهرست اختصارات: اختصار، عبارت کامل انگلیسی و معادل فارسی
\newglossarystyle{myAbbrlist}{%
  \renewenvironment{theglossary}{}{}
  \renewcommand*{\glsgroupskip}{\vskip 7mm}
  \renewcommand*{\glsgroupheading}[1]{%
    \begin{LTR}\subsection*{\glsgetgrouptitle{##1}}\end{LTR}}
  \renewcommand*{\glossentry}[2]{%
    \noindent
    \lr{\glsentrytext{##1}}\quad
    \lr{\Glsentrylong{##1}}\dotfill\space
    \glsentrydesc{##1}\par
  }
  \renewcommand*{\acronymname}{\rl{فهرست اختصارات}}
}

%%% تعریف دو واژه‌نامه
\newglossary[glg]{english}{gls}{glo}{واژه‌نامهٔ انگلیسی به فارسی}
\newglossary[blg]{persian}{bls}{blo}{واژه‌نامهٔ فارسی به انگلیسی}
\makeglossaries
\glsdisablehyper

%%% حفظ دستورات اصلی بسته پیش از بازتعریف
\let\oldgls\gls
\let\oldglspl\glspl

\makeatletter
\renewrobustcmd*{\gls}{\@ifstar\@msgls\@mgls}
\newcommand*{\@mgls}[1]{%
  \ifthenelse{\equal{\glsentrytype{#1}}{english}}
  {\oldgls{#1}\glsuseri{f-#1}}
  {\lr{\oldgls{#1}}}}
\newcommand*{\@msgls}[1]{%
  \ifthenelse{\equal{\glsentrytype{#1}}{english}}
  {\glstext{#1}\glsuseri{f-#1}}
  {\lr{\glsentryname{#1}}}}

\renewrobustcmd*{\glspl}{\@ifstar\@msglspl\@mglspl}
\newcommand*{\@mglspl}[1]{%
  \ifthenelse{\equal{\glsentrytype{#1}}{english}}
  {\oldglspl{#1}\glsuseri{f-#1}}
  {\oldglspl{#1}}}
\newcommand*{\@msglspl}[1]{%
  \ifthenelse{\equal{\glsentrytype{#1}}{english}}
  {\glsplural{#1}\glsuseri{f-#1}}
  {\glsentryplural{#1}}}
\makeatother

%%% تعریف هم‌زمان مدخل انگلیسی--فارسی و فارسی--انگلیسی
\newcommand{\newword}[4]{%
  \newglossaryentry{#1}{%
    type={english},
    name={\lr{#2}},
    plural={#4},
    text={#3},
    description={}}
  \newglossaryentry{f-#1}{%
    type={persian},
    name={#3},
    text={\lr{#2}},
    description={}}
}

%%% درج پاورقی انگلیسی در نخستین استفاده از واژه
\defglsentryfmt[english]{%
  \glsgenentryfmt
  \ifglsused{\glslabel}{}{\LTRfootnote{\glsentryname{\glslabel}}}}

%%% درج توضیح فارسی اختصار در نخستین استفاده
\defglsentryfmt[acronym]{%
  \glsentryname{\glslabel}
  \ifglsused{\glslabel}{}{\footnote{\glsentrydesc{\glslabel}}}}

%%% چاپ فهرست اختصارات
\let\Oldprintglossary\printglossary
\newcommand{\printabbreviation}{%
  {\renewcommand{\thepage}{\arabic{page}}\glsaddall[types={acronym}]}
  \begingroup
    \baselineskip=.75cm
    \setglossarystyle{myAbbrlist}
    \Oldprintglossary[type=acronym,title={فهرست اختصارات}]
  \endgroup
  \clearpage
}
\newcommand{\printacronyms}{\printabbreviation}

%%% چاپ هر دو واژه‌نامه؛ همهٔ مدخل‌های words.tex چاپ می‌شوند، حتی اگر در متن با \gls فراخوانی نشده باشند.
\renewcommand{\printglossary}{%
  \let\appendix\relax
  %\clearpage
  {\renewcommand{\thepage}{\arabic{page}}\glsaddall[types={persian,english}]}
  \twocolumn
  \setglossarystyle{myFaToEn}
  \Oldprintglossary[type=persian,title={واژه‌نامهٔ فارسی به انگلیسی},toctitle={واژه‌نامهٔ فارسی به انگلیسی}]
  \clearpage
  \setglossarystyle{myEntoFa}
  \Oldprintglossary[type=english,title={واژه‌نامهٔ انگلیسی به فارسی},toctitle={واژه‌نامهٔ انگلیسی به فارسی}]
  \onecolumn
}
